\documentclass{article}
\usepackage{tcee-repaso}
\tcetema[Política comercial estratégica]{3.B.8}{La política comercial (II). La política comercial estratégica.}
\tcefecha{05/11/2026}
\begin{document}
\cabecerarepaso
\begin{columnas}

% ==========================================================================
\bloque{Esquema del cante}
\esquemacante{%
\apartadocante{Introducción: «arrieros somos\ldots»}{2,5}
\apartadocante{I. Política comercial estratégica}{17}
\subcante{1. Idea}{2}
\subcante{2. Ofensiva: B-S, E-G, \emph{Airbus}-\emph{Boeing}}{7,5}
\subcante{3. Defensiva: aranceles, E-K, industria naciente}{5,5}
\subcante{4. Críticas}{2}
\apartadocante{II. Política de promoción exterior}{5}
\subcante{1. Objetivos}{1}
\subcante{2. Justificación}{2,5}
\subcante{3. Instrumentos}{1,5}
\apartadocante{III. Estrategia entre gobiernos}{4,5}
\subcante{1. Represalias en oligopolio}{1,5}
\subcante{2. Guerra arancelaria (\autor{Johnson}, 1953)}{2}
\subcante{3. Salida cooperativa}{1}
\apartadocante{Conclusión}{1}}
\emph{Reparto fino}: B-S \tiempo{4}, E-G \tiempo{1,5}, \emph{Airbus}-\emph{Boeing} \tiempo{2}; aranceles \tiempo{2,5}, E-K \tiempo{0,5}, ind. naciente \tiempo{2,5}. \emph{Pizarra}: tabla de modelos vacía en una esquina (en \hyperref[rep:III]{III}, columna «si responden») y las 2 reacciones recuadradas. Imprescindibles: reacciones con subvención y con arancel, costes medios de la industria naciente, matrices de \emph{Airbus}-\emph{Boeing}, tabla de represalias.

% ==========================================================================
\bloque[2,5]{Introducción}
\begin{itemize}
\item \textbf{Enganche}: «arrieros somos y en el camino nos encontraremos». UE y EE.\,UU. ante la OMC (2004) por las ayudas a \emph{Airbus} y \emph{Boeing}; represalias autorizadas $\approx$7.500 M\$/año a EE.\,UU. (2019) y $\approx$4.000 M\$ a la UE (2020); tregua de 5 años (2021), prorrogada por la UE (jul. 2026). ¿Puede un país ganar ayudando a sus empresas a competir fuera, si los demás pueden hacer lo mismo?
\item \textbf{Relevancia}: \emph{teórica}, primera grieta seria en el consenso librecambista desde el arancel óptimo; libre comercio = «regla práctica razonable», ya no lo óptimo (\autor{Krugman}, 1987: «ha perdido su inocencia»). \emph{Práctica}: CHIPS Act e IRA (EE.\,UU., 2022), Ley Europea de Chips (2023), derechos compensatorios a VE chinos (UE, 2024), aranceles de EE.\,UU. de 2025 \vertema{3.B.41}.
\item \textbf{Contextualización}: \concepto{política comercial} = medidas sobre volumen, composición o dirección del comercio exterior \vertema{3.B.7}. Comp. perfecta: país pequeño \ensuremath{\Rightarrow} intervenir ↓ bienestar; país grande \ensuremath{\Rightarrow} solo arancel óptimo (RRI). Nueva teoría (años 80) \vertema{3.B.6}: \concepto{traslado de rentas} en oligopolio (B-S, 1985) y creación de ventaja con escala (industria naciente).
\item \textbf{Problemática}: ¿se gana quitando beneficios a empresas extranjeras? ¿Es robusta la receta y qué debe saber el gobierno? ¿Qué pasa si los demás responden? ¿Qué alternativa compatible con la OMC queda?
\item \textbf{Estructura}: de la empresa al gobierno y a la interacción entre gobiernos: \hyperref[rep:I]{I}, un gobierno frente a empresas extranjeras (ofensiva, defensiva, críticas); \hyperref[rep:II]{II}, promoción exterior; \hyperref[rep:III]{III}, represalias y cooperación.
\end{itemize}

% ==========================================================================
\bloque[17]{I. Política comercial estratégica\label{rep:I}}
\begin{ideaclave}
Con competencia imperfecta, un gobierno que mueve primero puede trasladar rentas a su país (subvención con cantidades, impuesto con precios, arancel si la rival vende en casa) y, con economías externas, crear la ventaja que el accidente histórico impidió. Pero el instrumento depende de lo que el gobierno no sabe, la renta se disipa y todo descansa en que los demás no respondan.
\end{ideaclave}

\sub{I.1. Idea\label{rep:I.1}}
\begin{itemize}
\item \concepto{Política comercial estratégica}: la que usa el gobierno (subvenciones o impuestos a la exportación, aranceles) para cambiar a favor de sus empresas el resultado de su interacción estratégica con las extranjeras y ↑ bienestar nacional (\autor{Brander}, 1995). 3 elementos: \emph{quién} (el gobierno, mueve antes y se compromete), \emph{con qué}, \emph{para qué} (rentas del oligopolio o instalar una industria con escala).
\item Motor: \concepto{traslado de rentas} (\emph{profit shifting}). La distingue el \emph{motivo}, no el instrumento: no es la RRI, son las pocas empresas que se tienen en cuenta.
\item \emph{Justificación}: \concepto{segundo óptimo} (\autor{Lipsey}-\autor{Lancaster}, 1956) \vertema{3.A.22}: el oligopolio internacional es un fallo que el gobierno no puede eliminar. Contrapartida: \concepto{principio de focalización} (\autor{Bhagwati}-\autor{Ramaswami}, 1963): fallo interno \ensuremath{\Rightarrow} instrumento sobre él; el comercial es primer óptimo solo si el fallo está en el comercio.
\item Herramienta: juegos no cooperativos \vertema{3.A.14}, \autor{Cournot}/\autor{Bertrand} \vertema{3.A.19}. Juego en 2 etapas (gobierno \ensuremath{\Rightarrow} empresas), inducción hacia atrás \ensuremath{\Rightarrow} \concepto{ENPS}.
\item \concepto{Tercer mercado}: venden solo a un tercer país \ensuremath{\Rightarrow} $W_A$ = beneficio \ensuremath{-} coste de la política; traslado de rentas puro (ofensiva). \concepto{Mercados recíprocos}: cada una vende en los 2 países, segmentados (defensiva; dumping recíproco \vertema{3.B.6}).
\end{itemize}

\textbf{Tabla de modelos}: \concepto{ofensiva} (sobre exportaciones) y \concepto{defensiva} (sobre importaciones):
\begin{itemize}
\item \emph{B-S (1985)}, tercer mercado, \autor{Cournot} \ensuremath{\Rightarrow} \textbf{subvención} a la exportación; renta de la empresa extranjera.
\item \emph{E-G (1986)}, tercer mercado, \autor{Bertrand} \ensuremath{\Rightarrow} \textbf{impuesto} a la exportación; renta de los consumidores extranjeros.
\item \emph{Airbus-Boeing}, juego de entrada \ensuremath{\Rightarrow} subvención comprometida; la rival no entra.
\item \emph{B-S (1984)}, mercados recíprocos, \autor{Cournot} \ensuremath{\Rightarrow} \textbf{arancel}; renta de la extranjera.
\item \emph{E-K (1984)}, producto diferenciado \ensuremath{\Rightarrow} arancel; renta del monopolista extranjero.
\item \emph{Industria naciente}, comp. perfecta con economías externas \ensuremath{\Rightarrow} protección temporal (segundo óptimo); no se traslada renta: se crea la ventaja.
\end{itemize}

\sub{I.2. Ofensiva (1): \autor{Brander}-\autor{Spencer} (1985)\label{rep:I.2}}
\begin{itemize}
\item \emph{Idea}: oligopolio con beneficios extraordinarios protegidos por barreras (I+D aeronáutica). No se corrige el poder de mercado: se \emph{apropia} la renta.
\item \emph{Supuestos}: 2 países con una empresa cada uno, bien homogéneo \emph{solo para exportar}; demanda $p(Q)$, $Q=q_A+q_B$; coste marginal $c$; \autor{Cournot} (variación conjetural nula); gobierno de $A$ fija $s_A$ por unidad exportada antes y se compromete; $W_A$ = beneficio \ensuremath{-} subvención.
\item \emph{2.ª etapa}: CPO de $A$ y su función de reacción, que la subvención desplaza:
\funcionclave{p(Q)+p'(Q)\,q_A-\mathrm{CMg}_A+s_A=0\;\Rightarrow\;R_A(q_B;s_A)}
la de $B$, sin $s_A$, da $R_B(q_A)$.
\item \emph{Condiciones de \autor{Hahn}} (1962): $p'+q_ip''<0$ (cantidades = \concepto{sustitutos estratégicos}, $\partial^2B_i/\partial q_i\,\partial q_j<0$; \autor{Bulow}-\autor{Geanakoplos}-\autor{Klemperer}, 1985) y $p'-C_i''<0$ \ensuremath{\Rightarrow} $R_i'\in(-1,0)$: equilibrio de \autor{Cournot} existe, único (contracción) y estable.
\item \emph{1.ª etapa}: $W_A(s_A)=B_A-s_Aq_A=p(Q)\,q_A-C_A(q_A)$. Con la CPO:
\begin{formula}
$\displaystyle\frac{dW_A}{ds_A}=\underbrace{-s_A\frac{dq_A}{ds_A}}_{\text{producir de más}}+\underbrace{q_A\,p'(Q)\frac{dq_B}{ds_A}}_{\text{efecto estratégico}>0}$
\end{formula}
En $s_A=0$ solo queda el estratégico \ensuremath{\Rightarrow} algo de subvención siempre compensa.
\item \concepto{Subvención óptima}:
\begin{formula}
$s_A^*=q_A\,p'(Q)\,R_B'(q_A)>0$\quad(caso lineal: $s_A^*=\frac{a-c}{4}$)
\end{formula}
La CPO de $A$ pasa a ser la de una líder de \autor{Stackelberg}: $q_A=\frac{a-c}{2b}$, $q_B=\frac{a-c}{4b}$, $W_A=\frac{(a-c)^2}{8b}>\frac{(a-c)^2}{9b}$ (\autor{Cournot}). Con costes distintos, $s_A^*=(a-2c_A+c_B)/4$.
\item \emph{Gráfico}: $s_A$ desplaza $R_A$ a $R_A(q_B;s_A)$, que corta a $R_B$ en $S$, donde el isobeneficio de $A$ \emph{sin la subvención} (= $W_A$) es tangente a $R_B$.
\item \concepto{Compromiso creíble}: anunciar $q_A^S$ sin política no es creíble (volvería a $R_A$); la subvención, aprobada antes, sí (como la capacidad de \autor{Dixit}, 1980). «Poderoso caballero es don Dinero».
\end{itemize}
\graficorepaso[0.72\linewidth]{brander-spencer}{B-S (1985): la subvención lleva de $C$ al punto de \autor{Stackelberg} $S$}
\begin{itemize}
\item \emph{Implicaciones}: ENPS. Empresa nacional ↑ (incluso sin la subvención); extranjera ↓; consumidores del tercer país ↑ ($Q^S=\frac{3(a-c)}{4b}>Q^C=\frac{2(a-c)}{3b}$; «a caballo regalado\ldots»). \concepto{Empobrecimiento del vecino}; el mundo puede ganar. Descansa en 3 supuestos: gobierno mueve antes, sustitutos estratégicos, el otro gobierno no responde.
\item \emph{Extensiones}: \autor{Dixit} (1984), $n$ nacionales y $m$ extranjeras:
\begin{formula}
$s_A^*=\dfrac{(a-c)(m+1-n)}{2n(m+1)}>0\iff n<m+1$
\end{formula}
(si no, impuesto: ya compiten demasiado); libre entrada disipa la renta (\autor{Horstmann}-\autor{Markusen}, 1986); consumo nacional; I+D como compromiso (\autor{Spencer}-\autor{Brander}, 1983); entrada (\autor{Dixit}-\autor{Kyle}, 1985); información asimétrica; crecimiento \vertema{3.B.9}.
\item \emph{Evidencia}: simulaciones calibradas. \autor{Baldwin}-\autor{Krugman} (1988), A300 frente al 767: efecto sobre todo redistributivo; ganan los compradores. \autor{Krishna}-\autor{Hogan}-\autor{Swagel} (1989): ganancias pequeñas y de \emph{signo} cambiante. Acuerdo UE-EE.\,UU. de 1992 sobre ayudas \ensuremath{\Rightarrow} precios de los aviones +3,7\,\% (\autor{Irwin}-\autor{Pavcnik}, 2004).
\item \emph{Valoración}: con oligopolio una subvención puede ↑ $W$ (imposible en comp. perfecta). Límites: equilibrio parcial (con coste del fondo $1+\lambda$, $s^*$ menor), represalias, información.
\end{itemize}

\sub{I.2. Ofensiva (2): \autor{Eaton}-\autor{Grossman} (1986)}
\begin{itemize}
\item \emph{Idea}: \autor{Bertrand} (1883) objetó a \autor{Cournot} que se fijan precios; con solo ese cambio la receta se da la vuelta: \textbf{impuesto} a la exportación.
\item \emph{Supuestos}: los de B-S con competencia en \emph{precios} y \emph{bienes diferenciados} (evita la paradoja $p=c$) \vertema{3.A.18}; impuesto si $s_A<0$.
\item \emph{Desarrollo}: $R_A(p_B;s_A)$ y $R_B(p_A)$ \emph{crecientes}: precios = \concepto{complementos estratégicos}, $\partial^2B_i/\partial p_i\,\partial p_j>0$. $W_A=(p_A-c)\,q_A$ (antes del impuesto). Un impuesto ↑ $p_A$, la rival ↑ $p_B$ ($R_B'>0$) y ↑ ventas de $A$:
\begin{formula}
$s_A^*=(p_A-c)\,\dfrac{\partial q_A/\partial p_B}{\partial q_A/\partial p_A}\,R_B'(p_A)<0$
\end{formula}
misma regla que B-S con el signo de $R_B'$ cambiado. Lleva al liderazgo en precios de $A$.
\item \emph{Implicaciones}: país $A$ ↑ (su empresa puede perder neto, lo recauda el gobierno); empresa extranjera ↑; consumidores del tercer país ↓. Renta quitada a los \emph{consumidores extranjeros}: cártel impuesto por el gobierno.
\item \emph{Regla general}: conjetura frente a respuesta real. \autor{Cournot} \ensuremath{\Rightarrow} subvención; \autor{Bertrand} \ensuremath{\Rightarrow} impuesto; \concepto{conjeturas consistentes} \ensuremath{\Rightarrow} \emph{libre comercio}.
\item \emph{Valoración}: la política \emph{no es robusta}: un cambio menor de supuesto cambia subvención por impuesto, y el gobierno no observa cómo compiten. Es la crítica más seria.
\end{itemize}

\sub{I.2. Ofensiva (3): \emph{Airbus} frente a \emph{Boeing}}
\begin{itemize}
\item Versión discreta (\autor{Krugman}, 1987). Pagos (\emph{Boeing}, \emph{Airbus}): si entra una sola, 100; si entran las 2, $-5$ cada una; si no entra ninguna, 0.
\item Sin subvención: 2 equilibrios de \autor{Nash} (produce una). Si \emph{Boeing} se compromete antes: $(100,0)$.
\item Subvención europea de 10, comprometida \emph{antes}: producir es \emph{estrategia dominante} para \emph{Airbus} (5 si \emph{Boeing} entra, 110 si no) \ensuremath{\Rightarrow} \emph{Boeing} no entra: $(0,110)$; 100 de renta trasladada de EE.\,UU. a Europa. La subvención no paga producir: paga que la rival no entre.
\item Si \emph{Boeing} tuviera ventaja de costes (gana 5 aunque entren las 2): con la misma subvención entran las 2, \emph{Airbus} gana 5 con 10 de subvención y Europa pierde 5.
\end{itemize}

\sub{I.3. Defensiva (1): mercados recíprocos (\autor{Brander}-\autor{Spencer}, 1984)\label{rep:I.3}}
\begin{itemize}
\item \emph{Idea}: sin política, dumping recíproco (\autor{Brander}, 1981; \autor{Brander}-\autor{Krugman}, 1983) \vertema{3.B.6}. El arancel sustituye al coste de transporte: el dato pasa a ser instrumento.
\item \emph{Supuestos}: \concepto{mercados segmentados}, \autor{Cournot}, $p^A=a-bQ^A$, $Q^A=q_A^A+q_B^A$, sin transporte; arancel específico $t_A$ \ensuremath{\Rightarrow} coste de la extranjera en $A$: $c+t_A$. Basta estudiar el mercado $A$.
\item \emph{2.ª etapa}: funciones de reacción en el mercado $A$:
\funcionclave{R_A(q_B^A)=\tfrac{a-c}{2b}-\tfrac{1}{2}q_B^A,\quad R_B(q_A^A;t_A)=\tfrac{a-c-t_A}{2b}-\tfrac{1}{2}q_A^A}
$\partial q_A^A/\partial t_A=1/3b$; $\partial q_B^A/\partial t_A=-2/3b$; $\partial p^A/\partial t_A=1/3$.
\item \emph{1.ª etapa}: el gobierno maximiza
\begin{formula}
$W_A(t_A)=\underbrace{\tfrac{b}{2}(Q^A)^2}_{\mathrm{EC}^A}+\underbrace{(p^A-c)\,q_A^A}_{B_A^A}+\underbrace{t_A\,q_B^A}_{\text{recaudación}}$
\end{formula}
3 efectos: productores (nacional ↑), consumidores (↓ EC, deshace parte del efecto procompetitivo), recaudación (2/3 la paga la extranjera: \concepto{extracción de rentas}).
\item \concepto{Arancel estratégico}: $t_A^*=\frac{a-c}{3}>0$; $q_B^A=\frac{a-c}{9b}$, $q_A^A=\frac{4(a-c)}{9b}$; $\Delta W_A=\frac{(a-c)^2}{18b}$. Prohibitivo si $t_A\geq\frac{a-c}{2}$ (monopolio nacional).
\item \emph{Implicaciones}: el comercio no desaparece. El signo no está garantizado (demandas muy convexas \ensuremath{\Rightarrow} subvención a la importación). Si los 2 lo ponen \ensuremath{\Rightarrow} peor que el libre comercio.
\item \emph{Extensión}: \concepto{protección de las importaciones como promoción de las exportaciones} (\autor{Krugman}, 1984): CMg ↓ con la producción \emph{total}; reservar el mercado \ensuremath{\Rightarrow} nacional ↑ producción \ensuremath{\Rightarrow} ↓ CMg \ensuremath{\Rightarrow} ↑ ventas fuera (en cadena). Memorias de 16K: el mercado japonés protegido fue decisivo para exportar, pero costó a Japón más de lo que le dio (\autor{Baldwin}-\autor{Krugman}, 1986).
\item \emph{Evidencia}: aranceles de EE.\,UU. de 2018: los pagaron casi enteros los compradores estadounidenses (\autor{Amiti} et al., 2019; \autor{Fajgelbaum} et al., 2020) \ensuremath{\Rightarrow} apenas hubo extracción de rentas.
\item \emph{Valoración}: el motivo es la renta de la empresa, no la RRI (vale para un país pequeño). Límites: equilibrio parcial, CMg constante, segmentación, curvatura, represalias.
\end{itemize}
\graficorepaso[0.72\linewidth]{arancel-reciproco}{Mercados recíprocos: el arancel desplaza $R_B(q_A^A;t_A^*)$ hacia el origen; de $E_0$ a $E_1$}

\sub{I.3. Defensiva (2): \autor{Eaton}-\autor{Kierzkowski} (1984)}
\begin{itemize}
\item Producto diferenciado, mercados \emph{no} segmentados, demanda a la \autor{Lancaster} con \emph{precio crítico} \vertema{3.A.18}. $A$ no produce el bien; el monopolista de $B$ cobra en $A$ justo el precio crítico \ensuremath{\Rightarrow} $A$ igual que en autarquía.
\item Arancel \ensuremath{\Rightarrow} el precio a los consumidores no cambia (sigue el crítico): lo paga el monopolista de su renta. Óptimo: el que lo deja indiferente entre exportar o no; $W_A$ ↑ en la recaudación. Límite: caso muy particular.
\end{itemize}

\sub{I.3. Defensiva (3): la industria naciente}
\begin{itemize}
\item \emph{Idea}: industria que podría competir no se instala porque al principio sus costes son altos y nadie da el primer paso \ensuremath{\Rightarrow} protección \emph{temporal}. Argumento más antiguo: \autor{Hamilton} (1791), \autor{List} (1841), \autor{Mill} (única excepción), \autor{Chang} (2002); ISI en América Latina \vertema{3.B.28}.
\item Requisitos a la vez: \emph{rendimientos crecientes} (externos o aprendizaje), \emph{fallo de información o financiación} y \emph{accidente histórico}.
\item \emph{Modelo} (botones, \autor{Krugman}-\autor{Obstfeld}-\autor{Melitz}): \concepto{economías de escala externas} (\autor{Marshall}, 1890), comp. perfecta, $p=\mathrm{CMe}$. País 2 con $\mathrm{CMe}_2<\mathrm{CMe}_1$, pero la industria empezó en 1. Sin política, una empresa del 2 tendría coste $c_0>p_1$ \ensuremath{\Rightarrow} nadie entra. Con protección, la industria del 2 crece y baja por $\mathrm{CMe}_2$; al quitarla, el mundo queda en el punto 2 con $p_2<p_1$.
\item Por qué no lo hace el mercado: \emph{imperfecciones del mercado de capitales} y \emph{apropiabilidad} (los seguidores aprovechan lo que pagan los pioneros). En ambos casos la protección es un segundo óptimo.
\end{itemize}
\graficorepaso[0.78\linewidth]{industria-naciente}{Industria naciente: protección temporal y paso del punto 1 al 2}
\begin{itemize}
\item \emph{3 criterios} (\autor{Kemp}, 1960): \concepto{\autor{Mill}} (llegará a competir sin protección), \concepto{\autor{Bastable}} (VA de las ganancias > coste de la protección), \concepto{\autor{Kemp}} (las ganancias no las puede apropiar el pionero). El de \autor{Kemp} es el decisivo:
  \begin{itemize}
  \item economías \emph{internas} o aprendizaje propio \ensuremath{\Rightarrow} basta el mercado (salvo fallo de crédito: corregirlo allí);
  \item economías \emph{externas} o aprendizaje que se difunde \ensuremath{\Rightarrow} fallo de coordinación, equilibrios múltiples \ensuremath{\Rightarrow} protección o subvención justificable.
  \end{itemize}
\item Focalización: corregir el fallo en origen; el arancel, \concepto{segundo óptimo} (\autor{Baldwin}, 1969). «Muerto el perro, se acabó la rabia».
\item \emph{Evidencia}: mixta en la historia (raíles de acero, \autor{Head}, 1994; hojalata y arancel McKinley, \autor{Irwin}, 2000: habría llegado $\approx$10 años después); más favorable en experimentos naturales (bloqueo napoleónico, \autor{Juhász}, 2018; Corea 1973-1979, \autor{Lane}, 2025). Funcionan apoyos temporales y condicionados (\autor{Juhász}-\autor{Lane}-\autor{Rodrik}, 2024).
\item \emph{Valoración}: el único argumento que admitían los clásicos y el que mejor resiste (\autor{Melitz}, 2005: depende de cuánto y qué rápido se aprende). Límites: elegir la industria, búsqueda de rentas, lo temporal se vuelve permanente, represalias.
\end{itemize}

\sub{I.4. Críticas\label{rep:I.4}}
\begin{itemize}
\item \emph{Información}: el signo cambia con el modo de competencia, el número de empresas y la curvatura; el gobierno no conoce costes ni conducta; política del modelo equivocado \ensuremath{\Rightarrow} ↓ bienestar (K-H-S, 1989).
\item \emph{Disipación de la renta}: libre entrada (H-M, 1986); equilibrio parcial (impuestos distorsionantes, recursos de otros sectores).
\item \emph{Economía política}: búsqueda de rentas y captura \vertema{3.A.20}. «Quien paga, manda».
\item \emph{Represalias y reglas}: empobrecimiento del vecino \ensuremath{\Rightarrow} dilema del prisionero. Acuerdo SMC de la OMC: \emph{prohíbe} subvenciones supeditadas a la exportación o al contenido nacional (art.~3), recurribles las demás \vertema{3.B.33}.
\item \autor{Krugman} (1987): el libre comercio, mejor regla práctica, aunque no siempre lo óptimo. De ahí, la promoción exterior.
\end{itemize}

% ==========================================================================
\bloque[5]{II. Política de promoción exterior\label{rep:II}}
\begin{ideaclave}
Apoya la internacionalización, no cada exportación: corrige fallos de información, coordinación y financiación, baja el coste fijo de exportar y deja entrar a las empresas que podrían competir. No traslada rentas \ensuremath{\Rightarrow} no invita a la represalia y es compatible con la OMC.
\end{ideaclave}

\sub{II.1. Objetivos\label{rep:II.1}}
\begin{itemize}
\item \concepto{Defensiva}: ↓ importaciones. \concepto{Ofensiva}: sobre exportaciones; \emph{tradicional} (apoyo directo: financiación, seguro, subvenciones, ventajas fiscales) o \emph{moderna} (promoción exterior).
\item \concepto{Política de promoción exterior}: \emph{apoyo genérico} (no ligado al volumen exportado \ensuremath{\Rightarrow} no es subvención prohibida por el SMC) y directo a la \emph{internacionalización} (exportación + IED \vertema{3.B.18}; lazos estables, no éxito puntual).
\item \emph{Objetivo principal}: crecimiento y empleo. \emph{Intermedios}: competitividad y acceso a mercados (cadenas globales de valor).
\item España: política comercial cedida a la UE (1986) \vertema{3.B.41}; FIEM (Ley 11/2010), Ley 14/2013 de emprendedores.
\end{itemize}

\sub{II.2. Justificación\label{rep:II.2}}
\begin{itemize}
\item \emph{¿Por qué más internacionalización?} 5 motivos: \emph{productividad} (exportadoras más grandes y productivas ya antes, \autor{Bernard}-\autor{Jensen}, 1999; \autor{Melitz}, 2003; IED, umbral mayor, \autor{Helpman}-\autor{Melitz}-\autor{Yeaple}, 2004) \vertema{3.B.6}; \emph{innovación}; \emph{competencia} (↓ ineficiencia X, \autor{Leibenstein}, 1966); \emph{diversificación del riesgo}; \emph{externalidades} (autodescubrimiento, \autor{Hausmann}-\autor{Rodrik}, 2003).
\item \emph{¿Por qué intervenir?} Barreras = fallos de mercado: \emph{internas} (información con rasgos de \emph{bien público}; preparación del personal), \emph{externas} (encontrar clientes, proteccionismo, corrupción) y \emph{financiación} (\emph{información asimétrica}, pesa más con pymes).
\item En \autor{Melitz}: ↓ \emph{coste fijo de exportar} \ensuremath{\Rightarrow} ↓ umbral de productividad \ensuremath{\Rightarrow} entran las que estaban justo debajo. No toca el precio exportado \ensuremath{\Rightarrow} no traslada rentas.
\end{itemize}

\sub{II.3. Instrumentos\label{rep:II.3}}
\begin{itemize}
\item 5 grupos: \concepto{promoción comercial} (marca país, ferias, misiones, consorcios), \concepto{formación}, \concepto{información}, \concepto{diplomacia comercial} (la más sutil) y \concepto{financiación} (cofinanciación; crédito oficial solo si respeta el \concepto{Consenso de la OCDE}, anexo I-k del SMC; modernizado en 2023).
\item España: ICEX Next \vertema{4.A.25}; Estrategia PASE (12 países, 2017); Estrategia de Internacionalización 2017-2027; Plan de Respuesta y Relanzamiento Comercial (2025); FIEM, CESCE, COFIDES \vertema{4.A.24}.
\end{itemize}

% ==========================================================================
\bloque[4,5]{III. Estrategia entre gobiernos\label{rep:III}}
\begin{ideaclave}
Cuando los gobiernos responden, la política estratégica suele acabar en un dilema del prisionero (salvo los impuestos con competencia en precios, que perjudican al tercero). La salida: reciprocidad y no discriminación del GATT y reglas de la OMC sobre subvenciones: un juego repetido con reglas.
\end{ideaclave}

\sub{III.1. Represalias en los modelos de oligopolio\label{rep:III.1}}
\begin{itemize}
\item \emph{Subvenciones} (B-S): cada gobierno $s=\frac{a-c}{5}$ (algo menos que solo); cada país $W=\frac{2(a-c)^2}{25b}<\frac{(a-c)^2}{9b}$. Ganan las \emph{empresas} ($\frac{4(a-c)^2}{25b}$ cada una) y el tercer país; \emph{pierden los 2 países}.
\item \emph{Impuestos} (E-G): los 2 países \emph{ganan} (cártel); pierden el tercer país y el mundo. No es un dilema del prisionero.
\item \emph{Aranceles} (recíprocos): $t^*=\frac{a-c}{3}$ cada uno; $W=\frac{65(a-c)^2}{162b}<\frac{4(a-c)^2}{9b}$: Pareto-inferior al librecambio.
\item Con subvenciones y aranceles, intervenir es la mejor respuesta haga lo que haga el otro y los 2 acaban peor \ensuremath{\Rightarrow} \concepto{dilema del prisionero} \vertema{3.A.14}.
\end{itemize}

\sub{III.2. La guerra arancelaria (\autor{Johnson}, 1953)\label{rep:III.2}}
\begin{itemize}
\item Comp. perfecta, país grande: \concepto{arancel óptimo}
\begin{formula}
$t^{\mathrm{opt}}=\dfrac{1}{\varepsilon^*-1}$,\quad $\varepsilon^*>1$ elasticidad de la oferta neta del resto del mundo ($\varepsilon^*\to\infty$ \ensuremath{\Rightarrow} $t=0$)
\end{formula}
\item Sucesión de represalias con curvas de oferta neta y de indiferencia comercial (\autor{Meade}, 1952) \vertema{3.B.5}. Fin: \concepto{equilibrio de \autor{Nash} en aranceles} (cada curva de indiferencia comercial tangente a la oferta neta del otro); el comercio no desaparece.
\item Normalmente pierden los 2; pero \emph{puede ganar el país de oferta neta más elástica} (cuando la del otro es mucho menos elástica): en la práctica, \emph{el país mucho mayor} que su socio (\autor{Kennan}-\autor{Riezman}, 1988).
\item \emph{Evidencia}: fuera de la OMC, aranceles más altos donde la oferta extranjera es más inelástica (\autor{Broda}-\autor{Limão}-\autor{Weinstein}, 2008).
\end{itemize}
\graficorepaso[0.7\linewidth]{johnson}{Guerra arancelaria: equilibrio de \autor{Nash} en aranceles (caso simétrico: pierden los 2)}

\sub{III.3. La salida cooperativa\label{rep:III.3}}
\begin{itemize}
\item \autor{Bagwell}-\autor{Staiger} (1999): el GATT neutraliza el traslado de costes vía RRI con \emph{reciprocidad} y \emph{no discriminación} (NMF) \vertema{3.B.33}. Dilema del prisionero repetido: la amenaza de represalia sostiene la cooperación. «Dos no riñen si uno no quiere».
\item SMC: prohíbe subvenciones a la exportación y permite responder a las que dañan (\hyperref[rep:I.4]{I.4}).
\item \emph{Airbus-Boeing}: subvencionaron los 2, siguieron las 2 empresas y represalias a los 2 lados sobre vino, aceite, queso, tractores (DS316, DS353); tregua de 2021.
\item 2018-2019: aranceles de EE.\,UU. y respuestas; no apareció la mejora de la RRI. 2025: aranceles «recíprocos» con ley de emergencia; feb. 2026, el Supremo resuelve que esa ley no autoriza aranceles \vertema{3.B.41}.
\item \emph{Nueva carrera de subvenciones}: a la producción y la inversión (lo que las reglas permiten), por seguridad de suministro y externalidades; si todos subvencionan, las rentas van a consumidores y empresas y pagan los contribuyentes.
\end{itemize}

% ==========================================================================
\bloque[1]{Conclusión}
\begin{itemize}
\item \textbf{Recapitulación}: con competencia imperfecta se gana trasladando rentas (subvención con cantidades, impuesto con precios, arancel en casa) o creando ventaja con escala; pero el signo depende de lo no observable, la renta se disipa y, si responden, pierden los 2.
\item \textbf{Relevancia}: por eso la OMC prohíbe las subvenciones a la exportación y los desarrollados pasan a la promoción exterior; hoy, subvenciones industriales y aranceles a gran escala.
\item \textbf{Extensiones}: economía política de la protección (\autor{Grossman}-\autor{Helpman}, 1994); I+D y crecimiento \vertema{3.B.9}; política comercial común \vertema{3.B.41}; OMC \vertema{3.B.33}; instrumentos españoles \vertema{4.A.24}.
\item \textbf{Opinión}: mejor diagnóstico que receta: con externalidades claras, apoyo temporal, condicionado y evaluado; con solo rentas que trasladar, una carrera que pagan los contribuyentes.
\item \textbf{Idea final}: un país puede ganar una vez y mientras el otro no responda; «arrieros somos»: la ganancia duradera está en cooperar con reglas comunes.
\end{itemize}

% ==========================================================================
\bloque{Autores y fechas}
\hito{1791}{\autor{Hamilton}: protección de las manufacturas nacientes}
\hito{1841}{\autor{List}: industria naciente en Alemania}
\hito{1953}{\autor{Johnson}: guerra arancelaria}
\hito{1956}{\autor{Lipsey}-\autor{Lancaster}: segundo óptimo}
\hito{1960}{\autor{Kemp}: criterios de la industria naciente}
\hito{1962}{\autor{Hahn}: condiciones de estabilidad de \autor{Cournot}}
\hito{1963}{\autor{Bhagwati}-\autor{Ramaswami}: principio de focalización}
\hito{1983}{\autor{Brander}-\autor{Krugman}: dumping recíproco; \autor{Spencer}-\autor{Brander}: I+D}
\hito{1984}{\autor{Brander}-\autor{Spencer}: arancel; \autor{Krugman}: protección como promoción; \autor{Eaton}-\autor{Kierzkowski}; \autor{Dixit}: $n$ empresas}
\hito{1985}{\autor{Brander}-\autor{Spencer}: subvención a la exportación; \autor{Bulow}-\autor{Geanakoplos}-\autor{Klemperer}: sustitutos y complementos estratégicos}
\hito{1986}{\autor{Eaton}-\autor{Grossman}: impuesto; \autor{Horstmann}-\autor{Markusen}: libre entrada}
\hito{1987}{\autor{Krugman}: «¿ha pasado de moda el libre comercio?»; \emph{Airbus}-\emph{Boeing}}
\hito{1988}{\autor{Baldwin}-\autor{Krugman}: A300; \autor{Kennan}-\autor{Riezman}: gana el grande}
\hito{1999}{\autor{Bagwell}-\autor{Staiger}: economía del GATT}
\hito{2004}{\autor{Irwin}-\autor{Pavcnik}: acuerdo de aviones de 1992}
\hito{2018}{\autor{Juhász}: bloqueo napoleónico}

% ==========================================================================
\bloque{Datos clave}
\begin{itemize}
\item Represalias por \emph{Airbus}: \dato{$\approx$7.500 M\$/año}{OMC, 2019}; por \emph{Boeing}: \dato{$\approx$4.000 M\$}{OMC, 2020}; tregua \dato{5 años}{USTR, 2021}.
\item CHIPS Act: \dato{52.700 M\$}{Casa Blanca, 2022}.
\item Acuerdo de 1992 \ensuremath{\Rightarrow} precios de los aviones \dato{+3,7\,\%}{Irwin y Pavcnik, 2004}.
\item Agencias de promoción: \dato{$\approx$40 \$ exportados por \$}{Lederman et al., 2010}; \dato{$\approx$15 \$}{Olarreaga et al., 2015}.
\item Red exterior: \dato{107 Oficinas Económicas y Comerciales en 86 países}{ICEX, 2026}.
\item Pymes: \dato{99,8\,\% de las empresas, 61\,\% del empleo}{Ministerio de Industria, 2026}.
\item Exportadores regulares: \dato{+33,7\,\% en 2016 frente a 2011}{Secretaría de Estado de Comercio, 2017}.
\item ICEX Next: \dato{60\,\% de los gastos, máx. 24.000 € en 24 meses}{ICEX, 2026}.
\item Plan de Respuesta: \dato{14.100 M€}{La Moncloa, 2025}.
\item Consenso OCDE: \dato{hasta 15 años, 22 para inversiones verdes}{OCDE, 2023}.
\end{itemize}

% ==========================================================================
\bloque{Preguntas del tribunal}
\begin{itemize}
\item \textbf{¿Es la subvención de B-S un duopolio de \autor{Stackelberg}?} No: las empresas eligen a la vez; mueve primero el gobierno. $A$ está en $S$ \emph{sobre} su nueva reacción $R_A(q_B;s_A)$; la líder de \autor{Stackelberg} está fuera de $R_A$.
\item \textbf{¿Por qué E-G da un impuesto?} Precios = complementos estratégicos ($R_B'>0$): subir $p_A$ hace que la rival suba $p_B$. La política depende de conjetura frente a respuesta real.
\item \textbf{¿Cuándo es mejor un impuesto en B-S?} Con muchas nacionales, $n>m+1$ (\autor{Dixit}, 1984): compiten demasiado entre ellas.
\item \textbf{¿Por qué la promoción exterior no es subvención prohibida?} Apoyo genérico, no ligado al volumen exportado; baja el coste fijo, no el precio.
\item \textbf{¿Basta tener economías de escala para proteger?} No: con economías internas el pionero invierte solo (criterio de \autor{Kemp}); hace falta que no pueda apropiarse las ganancias.
\item \textbf{¿Quién gana la guerra de subvenciones?} Las empresas y el tercer país; pierden los 2 países.
\end{itemize}

% ==========================================================================
\begin{trampa}
\begin{itemize}
\item B-S: \emph{subvención} ($R_B'<0$, sustitutos estratégicos); E-G: \emph{impuesto} ($R_B'>0$, complementos); conjeturas consistentes: libre comercio.
\item En E-G el bienestar de $A$ sube aunque su empresa pueda ganar menos neto; la empresa \emph{extranjera} gana.
\item Guerra de subvenciones: ganan las \emph{empresas}, pierden los \emph{países}. Con impuestos (E-G), los 2 países \emph{ganan}.
\item Arancel recíproco: el precio sube solo 1/3 del arancel (los 2/3 los paga la extranjera); la extranjera pierde el doble de ventas de lo que gana la nacional ($-2/3b$ frente a $1/3b$).
\item \autor{Johnson}: el comercio no desaparece; puede ganar el país de oferta neta \emph{más} elástica (el mucho mayor), no siempre pierden los 2.
\item Criterios: \autor{Mill} = llegará a competir; \autor{Bastable} = VA ganancias > costes; \autor{Kemp} = no apropiables. Los 3 los recoge \autor{Kemp} (1960).
\item SMC: subvenciones a la exportación \emph{prohibidas} (art.~3); crédito oficial admitido si respeta el Consenso de la OCDE.
\end{itemize}
\end{trampa}

\end{columnas}
\end{document}
